\title{勉強ノート}
\author{}
\date{}

\usepackage{amsmath,amssymb}
\usepackage{amsthm}
\usepackage[OT2,T1]{fontenc}
\usepackage{mathtools}
\usepackage{mathrsfs}
\usepackage{otf}
\usepackage{url}
\usepackage[dvipdfmx]{hyperref}
\usepackage{pxjahyper}
\usepackage{enumitem}
\usepackage{fancyhdr}
\usepackage{tikz}
\usepackage[most]{tcolorbox}
\usetikzlibrary{cd,decorations.pathmorphing}

\DeclareMathOperator{\id}{id}
\DeclareMathOperator{\Hom}{Hom}
\DeclareMathOperator{\End}{End}
\DeclareMathOperator{\Ker}{Ker}
\DeclareMathOperator{\Coker}{Coker}
\DeclareMathOperator{\Image}{Im}
\DeclareMathOperator{\Spec}{Spec}
\DeclareMathOperator{\Frac}{Frac}
\DeclareMathOperator{\Fil}{Fil}
\DeclareMathOperator{\Rep}{Rep}
\DeclareMathOperator{\rank}{rank}
\DeclareMathOperator{\ord}{ord}
\DeclareMathOperator{\Aut}{Aut}
\DeclareMathOperator{\Gal}{Gal}
\DeclareMathOperator{\Tr}{Tr}
\DeclareMathOperator{\Norm}{N}
\DeclareMathOperator{\GL}{GL}
\DeclareMathOperator{\EssIm}{EssIm}

\DeclareSymbolFont{cyrletters}{OT2}{wncyr}{m}{n}
\DeclareMathSymbol{\Sha}{\mathalpha}{cyrletters}{"58}

\newcommand{\sourcepage}[1]{}
\newcommand{\noterefpage}[1]{}
\newcommand{\handref}[1]{}
\newcommand{\todo}[1]{\par\noindent\textbf{TODO.} #1\par}

\definecolor{noteLightLine}{gray}{0.82}
\definecolor{noteBadgeBack}{gray}{0.80}
\definecolor{noteLocalQed}{gray}{0.58}

\newcommand{\localqed}{\unskip\nobreak\mbox{\textcolor{noteLocalQed}{\(\blacklozenge\)}}}

\makeatletter
\newcommand{\noteenvtitle}[2]{%
    \def\note@arg{#2}%
    \ifx\note@arg\@empty #1\else #2\fi}
\newcommand{\noteenvhead}[2]{%
    \par\medskip\noindent\textbf{\noteenvtitle{#1}{#2}.}\ }
\newcommand{\noteauxlabel}[2]{%
    #1%
    \def\note@arg{#2}%
    \ifx\note@arg\@empty\else\ (#2)\fi}
\newcommand{\noteauxstart}[2]{%
    \par\medskip
    \begin{list}{}{\leftmargin=1.5em\rightmargin=0pt\topsep=0pt\parsep=0.3\baselineskip\itemsep=0pt}%
    \item[]\textbf{\noteauxlabel{#1}{#2}.}\ }
\newcommand{\noteauxend}{\end{list}\medskip}
\newcommand{\notelabel}[1]{\phantomsection\label{#1}}
\newcommand{\dfnref}[1]{\hyperref[dfn:#1]{定義~#1}}
\newcommand{\propref}[1]{\hyperref[prop:#1]{命題~#1}}
\newcommand{\thmref}[1]{\hyperref[thm:#1]{定理~#1}}
\newcommand{\corref}[1]{\hyperref[cor:#1]{系~#1}}
\newcommand{\lemref}[1]{\hyperref[lem:#1]{補題~#1}}
\newcommand{\claimref}[1]{\hyperref[claim:#1]{主張~#1}}
\newcommand{\rmkref}[1]{\hyperref[rmk:#1]{注意~#1}}
\newcommand{\exref}[1]{\hyperref[ex:#1]{例~#1}}
\newcommand{\exerciseref}[1]{\hyperref[exercise:#1]{演習~#1}}
\newcommand{\namedref}[2]{\hyperref[#1]{#2}}
\newcommand{\noteNextSectionTarget}{}
\newcommand{\noteNextChapterTarget}{}
\newcommand{\notepagelinks}[2]{%
    \gdef\noteNextSectionTarget{#1}%
    \gdef\noteNextChapterTarget{#2}}
\newcommand{\notenavsep}{\hspace{0.5em}{\normalfont |}\hspace{0.5em}}
\newcommand{\notenavlink}[2]{%
    \edef\note@target{#1}%
    \ifx\note@target\@empty
        \textcolor{gray}{#2}%
    \else
        \hyperref[\note@target]{#2}%
    \fi}
\newcommand{\noteifnonempty}[2]{%
    \edef\note@target{#1}%
    \ifx\note@target\@empty\else #2\fi}
\newcommand{\noteheaderlinks}{%
    \footnotesize\sffamily
    \hyperref[toc]{Contents}%
    \noteifnonempty{\noteNextSectionTarget}{\notenavsep\notenavlink{\noteNextSectionTarget}{Next subsubsection}}%
    \noteifnonempty{\noteNextChapterTarget}{\notenavsep\notenavlink{\noteNextChapterTarget}{Next chapter}}}
\makeatother

\ExplSyntaxOn
\cs_new_protected:Npn \note_print_title_tail:nn #1#2
  {
    \tl_set:Nn \l_tmpa_tl {#2}
    \tl_if_blank:VTF \l_tmpa_tl
      {}
      {
        \regex_replace_once:nnN { \A \s* #1 \s* } {} \l_tmpa_tl
        \tl_if_blank:VTF \l_tmpa_tl {} { \hspace{0.45em}\textbf{\tl_use:N \l_tmpa_tl} }
      }
  }
\NewDocumentCommand{\noteenvtail}{mm}{\note_print_title_tail:nn{#1}{#2}}
\ExplSyntaxOff

\newcommand{\noteenvbadge}[1]{%
    \tcbox[
        on line,
        box align=base,
        boxrule=0pt,
        arc=2pt,
        left=0.3em,
        right=0.3em,
        top=0.08em,
        bottom=0.08em,
        colback=noteBadgeBack,
        colframe=white
    ]{\bfseries #1}}
\newcommand{\noteinlineenvhead}[2]{%
    \par\medskip\noindent\noteenvbadge{#1\noteenvtail{#1}{#2}}\quad\ignorespaces}
\newcommand{\notelineenvhead}[2]{%
    \noindent\textbf{#1\noteenvtail{#1}{#2}.}\quad\ignorespaces}

\newtcolorbox{notelineenvbox}[1]{%
    enhanced jigsaw,
    breakable,
    frame hidden,
    interior hidden,
    boxrule=0pt,
    sharp corners,
    borderline west={4pt}{0pt}{#1},
    left=1.65em,
    right=0pt,
    top=0.2em,
    bottom=0.2em,
    before skip=0.9\baselineskip,
    after skip=0.9\baselineskip}

\newenvironment{dfn}[1][]{\noteinlineenvhead{定義}{#1}}{\par\medskip}
\newenvironment{prop}[1][]{\begin{notelineenvbox}{noteLightLine}\notelineenvhead{命題}{#1}}{\end{notelineenvbox}}
\newenvironment{thm}[1][]{\begin{notelineenvbox}{black}\notelineenvhead{定理}{#1}}{\end{notelineenvbox}}
\newenvironment{cor}[1][]{\begin{notelineenvbox}{noteLightLine}\notelineenvhead{系}{#1}}{\end{notelineenvbox}}
\newenvironment{lem}[1][]{\begin{notelineenvbox}{noteLightLine}\notelineenvhead{補題}{#1}}{\end{notelineenvbox}}
\newenvironment{claim}[1][]{\begin{notelineenvbox}{noteLightLine}\notelineenvhead{主張}{#1}}{\end{notelineenvbox}}
\newenvironment{rmk}[1][]{\noteenvhead{注意}{#1}}{\par\medskip}
\newenvironment{ex}[1][]{\noteenvhead{例}{#1}}{\par\medskip}
\newenvironment{question}[1][]{\noteenvhead{問}{#1}}{\par\medskip}
\newenvironment{conjecture}[1][]{\noteenvhead{予想}{#1}}{\par\medskip}
\newenvironment{recall}[1][]{\noteauxstart{確認}{#1}}{\noteauxend}
\newenvironment{localdfn}[1][]{\noteauxstart{定義}{#1}}{\noteauxend}
\newenvironment{gapnote}{\par\medskip\noindent\textbf{行間補足.}\ }{\par\medskip}
\newenvironment{finalnote}{\par\medskip\noindent\ignorespaces}{\par\medskip}

\setlength{\headheight}{14pt}
\pagestyle{fancy}
\fancyhf{}
\fancyhead[R]{\noteheaderlinks}
\fancyfoot[C]{\thepage}
\renewcommand{\headrulewidth}{0pt}
\renewcommand{\footrulewidth}{0pt}
\fancypagestyle{plain}{%
    \fancyhf{}%
    \fancyhead[R]{\noteheaderlinks}%
    \fancyfoot[C]{\thepage}%
    \renewcommand{\headrulewidth}{0pt}%
    \renewcommand{\footrulewidth}{0pt}}

\makeatletter
\def\@seccntformat#1{\csname the#1\endcsname\ }
\renewenvironment{proof}[1][\proofname]{\par
    \pushQED{\qed}%
    \normalfont \topsep6\p@\@plus6\p@\relax
    \trivlist
    \item\relax
    {\bfseries
    #1\@addpunct{.}}\hspace\labelsep\ignorespaces
}{%
    \popQED\endtrivlist\@endpefalse
}

\renewcommand{\proofname}{証明}
\makeatother
